% Figure: muscle force-length and force-velocity curves (Hill model).
% Two panels side by side.
\begin{figure}[htbp]
\centering
\begin{tikzpicture}[
  axis/.style={-{Stealth}, thick, draw=engblack},
  curve/.style={very thick},
  annot/.style={font=\scriptsize, align=center},
]

% ===== Left panel: force-length =====
\begin{scope}[xshift=0cm]
\draw[axis] (0, 0) -- (6.2, 0) node[right, font=\small] {$L/L_0$};
\draw[axis] (0, 0) -- (0, 4.7) node[above, font=\small] {$F/F_0$};

\foreach \x/\xt in {1/0.5, 2/0.75, 3/1.0, 4/1.25, 5/1.5} {
  \draw (\x, 0) -- (\x, -0.08) node[below, font=\scriptsize] {\xt};
}
\foreach \y/\yt in {1/0.25, 2/0.5, 3/0.75, 4/1.0} {
  \draw (0, \y) -- (-0.08, \y) node[left, font=\scriptsize] {\yt};
}

% Active force-length: triangular about L/L0 = 1 with plateau
\draw[curve, engred, smooth] (1.0, 0) .. controls (2.0, 1.5) and (2.7, 3.5) .. (3.0, 4.0)
   -- (3.4, 4.0) .. controls (3.7, 4.0) and (4.5, 1.0) .. (5.0, 0);
\node[engred, font=\small, anchor=south] at (3.2, 4.0) {active};

% Passive force-length: rising exponentially above L/L0 ~ 1
\draw[curve, engblue, smooth, samples=50, domain=2.7:5.2]
  plot (\x, {0.05 * exp(1.0 * (\x - 2.7))});
\node[engblue, font=\small, anchor=west] at (5.0, 3.2) {passive};

% Total = active + passive (dashed)
\draw[curve, engblack, dashed, smooth, samples=50, domain=1.0:5.2]
  plot (\x, { (\x < 3.0 ? \x - 1.0 : (\x < 3.4 ? 4.0 : 4.0 - 8.0*(\x - 3.4)/(5.0 - 3.4))) + 0.05*exp(1.0*max(\x-2.7,0))});

\node[font=\bfseries, anchor=north] at (3.0, -0.7) {Force-length};
\end{scope}

% ===== Right panel: force-velocity =====
\begin{scope}[xshift=8.5cm]
\draw[axis] (0, 0) -- (6.2, 0) node[right, font=\small] {$v / v_{\max}$};
\draw[axis] (0, 0) -- (0, 4.7) node[above, font=\small] {$F / F_0$};

\foreach \x/\xt in {1/-0.5, 2/0, 3/0.5, 4/1.0} {
  \draw (\x+1.0, 0) -- (\x+1.0, -0.08) node[below, font=\scriptsize] {\xt};
}
\foreach \y/\yt in {1/0.5, 2/1.0, 3/1.5, 4/1.8} {
  \draw (0, \y) -- (-0.08, \y) node[left, font=\scriptsize] {\yt};
}

% Hill curve: concentric (v > 0): (F + a)(v + b) = (F_0 + a) b
% Take a/F0 = 0.25, b/vmax = 0.25
% F/F0 = ((F0+a) b / (v+b) - a) / F0
% Eccentric (v < 0): rises to 1.5-1.8 F0 then plateau
\draw[curve, engred, smooth, samples=80, domain=0:1.0]
  plot ({2.0 + 4.0*\x}, {2.0 * (1.25 * 0.25 / (\x + 0.25) - 0.25)});

% Eccentric branch
\draw[curve, engred, smooth, samples=40, domain=-0.5:0]
  plot ({2.0 + 4.0*\x}, {2.0 + 1.6 * (1 - exp(5*\x))});

% Mark isometric F_0 at v = 0
\fill[engred] (2.0, 2.0) circle (1.5pt);
\node[engred, anchor=south west, font=\scriptsize] at (2.1, 2.05) {$F_0$};
% Mark v_max at F = 0
\fill[engred] (6.0, 0) circle (1.5pt);
\node[engred, anchor=south, font=\scriptsize] at (6.0, 0.1) {$v_{\max}$};

\node[engred, font=\small, anchor=west] at (4.0, 0.9) {concentric};
\node[engred, font=\small, anchor=east] at (1.7, 3.4) {eccentric};

\node[font=\bfseries, anchor=north] at (3.0, -0.7) {Force-velocity (Hill)};
\end{scope}

\end{tikzpicture}
\caption{Hill-type muscle model: force-length (left) and force-velocity
(right) characteristics. Active force peaks at the optimal length
$L_0$ where actin-myosin overlap is maximal; passive force grows
exponentially above resting length as the connective-tissue elements
stretch. Concentric shortening at velocity $v$ produces force
$F(v) = (F_0 + a)\,b / (v + b) - a$, the Hill hyperbola, with
$a/F_0 \approx 0.25$ and $b/v_{\max} \approx 0.25$ for mammalian
skeletal muscle. Eccentric (lengthening) contraction produces forces
up to $\sim 1.5\,F_0$ to $1.8\,F_0$, the mechanical mechanism by
which downhill walking and the deceleration phase of running impose
the largest forces on tendons.}
\label{fig:vol06ch06:muscle-curves}
\end{figure}
